\documentclass[11pt]{article}
\usepackage[margin=1in]{geometry}
\usepackage{amsmath,amssymb,amsthm}
\usepackage{xurl}
\usepackage{hyperref}
\sloppy
\let\oldus\_ \renewcommand{\_}{\oldus\allowbreak}
\newtheorem{theorem}{Theorem}
\newtheorem{lemma}[theorem]{Lemma}
\theoremstyle{definition}
\newtheorem{definition}[theorem]{Definition}
\newtheorem{remark}[theorem]{Remark}

\title{Disproof of conjecture 00000000421:\\ the expected number of addable cells is at least $7/3$ at every prime size}
\author{Jackmeson1}
\date{2026-10-05}

\begin{document}
\maketitle

\section{The conjecture and the reading}

\textbf{Statement (English, verbatim).} Definition: The growth-tree statistic $a(T)$ is the number
of outer-corner cells that can be added to the tableau $T$. Conjecture: The expectation of $a(T)$
on a uniformly random tableau converges to two; and the convergence rate is the inverse logarithm
of the tableau size. (growth-tree expectation limit two) (The Chinese text states the same.)

\textbf{Reading.} The growth tree is Young's lattice: a tableau grows by adding one cell at a time,
and the cells that can be added to $T$ are the addable (outer-corner) cells of its shape. For
$n\ge 1$ let $E_n$ be the expectation of $a(T)$ when $T$ is a uniformly random standard Young
tableau with $n$ cells. The conjecture asserts
\[
  \text{(L)}\quad E_n \to 2 \ (n\to\infty)
  \qquad\text{and}\qquad
  \text{(R)}\quad |E_n-2| \text{ is of order } 1/\log n .
\]
Every version of (R) (namely $|E_n-2|=O(1/\log n)$, $\Theta(1/\log n)$, or $E_n-2\sim c/\log n$)
implies $|E_n-2|\le C/\log n$ for some constant $C$ and all large $n$. We refute (L), and we refute
this weakest form of (R). We also refute both clauses for a second reading of ``tableau'', namely a
Young diagram (Ferrers shape) with $n$ cells chosen uniformly.

\textbf{Readings not covered.} (i) Semistandard tableaux: with unbounded entries there are infinitely
many of each shape, so there is no uniform measure, and we do not treat bounded-entry variants.
(ii) Counting removable (inner) corners instead of addable ones: the text says cells that ``can be
added'', so we do not treat this. (iii) The Plancherel measure $f_\lambda^2/n!$: the counting
argument below applies to it as well (the row and the column have weight $1/n!$ each), but this is
not formalized.

\section{Definitions (as in the Lean file)}

\begin{definition}
A \emph{Young diagram} is Mathlib's \texttt{YoungDiagram}: a finite set $\mu\subset\mathbb N\times\mathbb N$
that is a lower set for the product order. Cell $(i,j)$ lies in row $i$ and column $j$, counted from
$0$. We write $r=\mathtt{rowLen}\,0$ and $c=\mathtt{colLen}\,0$ for the lengths of the first row and the
first column. Mathlib proves $(i,j)\in\mu \iff j<\mathtt{rowLen}\,i \iff i<\mathtt{colLen}\,j$.
\end{definition}

\begin{definition}[addable cells, \texttt{addable}, \texttt{a}]
A cell $x\notin\mu$ is \emph{addable} (an outer corner) if $\mu\cup\{x\}$ is again a lower set,
that is, a Young diagram. The statistic is $a(\mu)=\#\{\text{addable cells of }\mu\}$, written in Lean
as \texttt{Set.ncard}. We prove that this set is finite (\texttt{addable\_finite}).
\end{definition}

\begin{definition}[standard Young tableaux, \texttt{SYT}]
A standard Young tableau with $n$ cells is recorded by the injective map
$\mathrm{pos}:\{0,\dots,n-1\}\to\mathbb N\times\mathbb N$, where $\mathrm{pos}(k)$ is the cell holding the
entry $k+1$. Its image is required to be a lower set (the shape \texttt{SYT.shape}, a Young diagram
with $n$ cells), and entries must increase from left to right along each row and from top to
bottom down each column: if $\mathrm{pos}(k)$ and $\mathrm{pos}(l)$ are in the same row (column) and
$\mathrm{pos}(k)$ is to the left of (above) $\mathrm{pos}(l)$, then $k<l$. We prove that there are
finitely many such tableaux, since all coordinates are $<n$ (\texttt{SYT.pos\_lt}). The statistic of a
tableau is $a$ of its shape.
\end{definition}

\begin{definition}[the two expectations]
$E_{\mathrm{SYT}}(n)=\frac{1}{\#\mathrm{SYT}(n)}\sum_{T\in\mathrm{SYT}(n)} a(\mathrm{shape}\,T)$ (\texttt{E\_SYT}), and
$E_{\mathrm{shape}}(n)=\frac{1}{\#\mathrm{Shape}(n)}\sum_{\mu\in\mathrm{Shape}(n)} a(\mu)$ (\texttt{E\_shape}),
where $\mathrm{Shape}(n)$ is the set of Young diagrams with $n$ cells.
\end{definition}

\section{Proof}

\begin{lemma}[\texttt{corner\_row}, \texttt{corner\_col}, \texttt{two\_le\_a}]
The cells $(0,r)$ and $(c,0)$ are addable. If $\mu\ne\emptyset$, they are distinct, so $a(\mu)\ge 2$.
\end{lemma}
\begin{proof}
A cell $(i,j)\notin\mu$ is addable as soon as $(i-1,j)\in\mu$ when $i>0$ and $(i,j-1)\in\mu$ when
$j>0$ (\texttt{mem\_addable\_of}). Indeed, any $y<(i,j)$ in the product order is below one of these
two cells, so it lies in $\mu$. Now $(0,r)\notin\mu$, and if $r>0$ then $(0,r-1)\in\mu$; if $r=0$
the condition is vacuous (for the empty diagram $(0,r)=(0,0)$ is addable). The case of $(c,0)$ is
symmetric, using $c>0$ in the same way. If $(0,0)\in\mu$, then $r,c\ge 1$, so $(0,r)\ne(c,0)$.
\end{proof}

\begin{lemma}[\texttt{three\_le\_a}]
If some cell $(i,j)$ with $i<c$ and $j<r$ is not in $\mu$ (so $\mu$ is not the $c\times r$ rectangle),
then $a(\mu)\ge3$.
\end{lemma}
\begin{proof}
Let $k$ be the least row with $\mathtt{rowLen}\,k<r$. Such a row exists, since $\mathtt{rowLen}\,i\le j<r$, and
$k\le i<c$. Also $k\ge 1$, because row $0$ has length $r$. Since $k<c$ we have $(k,0)\in\mu$, so
$\mathtt{rowLen}\,k\ge1$. The cell $(k,\mathtt{rowLen}\,k)$ is not in $\mu$, and the cell to its left is in
$\mu$. The cell above it, $(k-1,\mathtt{rowLen}\,k)$, is also in $\mu$, because $\mathtt{rowLen}(k-1)\ge r>\mathtt{rowLen}\,k$
by minimality of $k$. So it is addable. It differs from $(0,r)$ because $k\ge1$, and from $(c,0)$
because $k<c$.
\end{proof}

\begin{lemma}[\texttt{rowOrCol\_or\_three\_le}]
If $\#\mu=p$ is prime, then $\mu$ is a single row ($\forall x\in\mu,\ x_1=0$), or a single column
($\forall x\in\mu,\ x_2=0$), or $a(\mu)\ge3$.
\end{lemma}
\begin{proof}
If the hypothesis of the previous lemma fails, every cell of the box $[0,c)\times[0,r)$ is in $\mu$.
Conversely, each cell $(i,j)\in\mu$ has $i<c$ and $j<r$. So $\mu$ is that box and $p=c\cdot r$. Since $p$
is prime, $c=1$ (a row) or $r=1$ (a column).
\end{proof}

\begin{lemma}[\texttt{SYT.row\_unique}, \texttt{SYT.col\_unique}, \texttt{shape\_eq\_of\_row}, \texttt{shape\_eq\_of\_col}]
There is at most one standard tableau with $n$ cells whose shape is a single row, and at most one
whose shape is a single column. Two single-row (or two single-column) diagrams with the same number
of cells are equal.
\end{lemma}
\begin{proof}
For a single-row tableau, the row condition and injectivity make $k\mapsto \mathrm{pos}(k)_2$ a strictly
increasing map $\{0,\dots,n-1\}\to\{0,\dots,n-1\}$. Mathlib's \texttt{Finset.orderEmbOfFin\_unique} shows
that such a map is unique, so any two single-row tableaux coincide. Columns are handled in the same
way. A single-row diagram with $n$ cells is $\{(0,j):j<n\}$, and the same holds for a single column.
\end{proof}

\begin{lemma}[\texttt{hook}, \texttt{hook\_not\_rowOrCol}]
For $n\ge3$, placing $1$ at $(0,0)$, $2$ at $(1,0)$, and $3,\dots,n$ at $(0,1),\dots,(0,n-2)$ gives a
standard tableau. Its shape (the hook $(n-1,1)$) is neither a row nor a column.
\end{lemma}

\begin{lemma}[\texttt{avg\_bound}]
Let $X$ be a finite set with a map $\mathrm{sh}:X\to\{\text{Young diagrams}\}$. Assume that every
$\mathrm{sh}(x)$ has a prime number of cells, that at most one $x$ has single-row shape and at most one
has single-column shape, and that some $x_0$ has another shape. Then
$\frac{1}{\#X}\sum_{x}a(\mathrm{sh}\,x)\ge 7/3$.
\end{lemma}
\begin{proof}
Let $B$ be the set of $x$ with row or column shape, so $|B|\le2$ and $|B|+1\le t:=\#X$. Lemma 1
(the shapes are nonempty) and Lemma 3 give $a\ge 3-[x\in B]$ pointwise, so $\sum_x a\ge 3t-|B|$. Since
$3|B|\le 2|B|+2\le 2t$, we get $3t-|B|\ge \tfrac73 t$.
\end{proof}

\begin{theorem}[\texttt{E\_SYT\_ge}, \texttt{E\_shape\_ge}]
For every prime $p\ge3$, $E_{\mathrm{SYT}}(p)\ge 7/3$ and $E_{\mathrm{shape}}(p)\ge7/3$.
\end{theorem}
\begin{proof}
Apply the previous lemma with $X=\mathrm{SYT}(p)$ and $\mathrm{sh}=\mathrm{shape}$, and with
$X=\mathrm{Shape}(p)$ and $\mathrm{sh}=\mathrm{id}$, taking $x_0$ to be the hook tableau or its shape.
\end{proof}

\begin{theorem}[\texttt{not\_limit\_two}, \texttt{disproof}]
Let $E$ be $E_{\mathrm{SYT}}$ or $E_{\mathrm{shape}}$. Then $E(n)\not\to2$. Moreover, there is no
constant $C$ such that $|E(n)-2|\le C/\log n$ for all large $n$.
\end{theorem}
\begin{proof}
Either statement would give $E(n)<7/3$ for all large $n$. For (R), this uses $C/\log n\to0$. But
there are arbitrarily large primes (\texttt{Nat.exists\_infinite\_primes}), and at each of them
$E\ge 7/3$.
\end{proof}

\begin{remark}[not used, not formalized]
In fact $E_{\mathrm{SYT}}(n)=t_{n+1}/t_n$, where $t_n$ is the number of standard tableaux with $n$
cells: deleting the largest entry is a bijection from tableaux with $n+1$ cells onto pairs (tableau
with $n$ cells, addable cell). OEIS A000085 (\url{https://oeis.org/A000085}) identifies $t_n$:
``Number of self-inverse permutations on n letters, also known as involutions; number of standard
Young tableaux with n cells.'' It also gives the recurrence ``a(n) = a(n-1) + (n-1)*a(n-2)''. Hence
$E_{\mathrm{SYT}}(n)\to\infty$. Our script checks $E_{\mathrm{SYT}}(n)=t_{n+1}/t_n$ numerically for
$n\le22$.
\end{remark}

\section{Lean formalization and axiom audit}

The file \texttt{lean/Conjecture421/Basic.lean} (Lean 4, Mathlib v4.33.1; its only import is
\texttt{Mathlib}) proves
\begin{verbatim}
theorem C421.disproof :
    (forall p : Nat, p.Prime -> 3 <= p -> 7/3 <= E_SYT p /\ 7/3 <= E_shape p) /\
    (not Tendsto E_SYT atTop (nhds 2) /\
      not exists C : Real,
        eventually (n in atTop), |E_SYT n - 2| <= C / Real.log n) /\
    (not Tendsto E_shape atTop (nhds 2) /\
      not exists C : Real,
        eventually (n in atTop), |E_shape n - 2| <= C / Real.log n)
\end{verbatim}
(written here in ASCII; the source uses the usual Unicode notation). All the lemmas above appear
under the names given. It uses Mathlib's \texttt{YoungDiagram} API (\texttt{rowLen}, \texttt{colLen},
\texttt{mem\_iff\_lt\_rowLen}, \texttt{mem\_iff\_lt\_colLen}, \texttt{rowLen\_eq\_card},
\texttt{colLen\_eq\_card}, \texttt{up\_left\_mem}). No code was taken from other submissions.
\texttt{lake build} succeeds, and \texttt{\#print axioms C421.disproof} reports only
\texttt{[propext, Classical.choice, Quot.sound]}. The file contains no \texttt{sorry}, no
\texttt{native\_decide} and no added axioms.

\end{document}
